% =============================================================================
%  SECTION 5 — C3: Architectural vs. Representational Approaches
%  Slides 5.1–5.6  |  Target time: ~10 minutes
%
%  Corresponds to dissertation Ch. 5 (Contribution 3)
%
%  Drop this \section{} call and the six frames that follow into main.tex
%  after the §4 C2 slides (slide 4.8), before §6 (Synthesis).
% =============================================================================

%%==========================================================================
%% CONTRIBUTION 3
%%==========================================================================

\begin{frame}{Two Varieties of Invariance}
    C2 showed that a \emph{representational} trick (homogenized observations
    with zero-masking) matches HAPPO at a fraction of the storage.
    
    \vspace{0.4em}
    \textbf{The Permutation Invariant Critic}
    (PIC)~\cite{liu2020b} replaces the standard centralized critic
    with a \textbf{graph neural network}:
    \begin{itemize}
        \item Each agent is a \emph{node} for the critic
        \item Agent heterogeneity is encoded via explicit
        \emph{one-hot attribute vectors} on each node.
        \item The output is a team-level value estimate that is
        invariant to agent ordering and flexible to team size.
    \end{itemize}
\end{frame}


\begin{frame}{Two Varieties of Invariance}

    \vspace{0.3em}
    \begin{columns}
        \begin{column}{0.48\textwidth}
            \begin{block}{Implicit Indication (Representational)}
                Invariance lives in the \emph{input encoding}.\\
                Standard MLP; identity inferred from mask pattern.
            \end{block}
        \end{column}
        \begin{column}{0.48\textwidth}
            \begin{block}{PIC (Architectural)}
                Invariance achieved through GNN convolution. \\
                With agents as nodes in the critic, 
                and observation elements as nodes in the agent.
            \end{block}
        \end{column}
    \end{columns}
    \vspace{1em}

    \small \textbf{Experimental control:} matched network capacities,
    identical hyperparameters, same HyperGrid environment;
    only the invariance mechanism differs.

    \note{About 2 minutes.}
\end{frame}


\begin{frame}{C3 Research Questions}
    \begin{itemize}
        \item[\textbf{RQ3.1}]
        To what extent do graph-based policies (PIC) improve learning
        efficiency vs.\ representational alternatives?
        \item[\textbf{RQ3.2}]
        How robust are graph-based policies to sensor configuration
        and team composition changes vs.\ representational approaches?
        \item[\textbf{RQ3.3}]
        What are the computational and implementation costs of
        architectural vs.\ representational invariance?
    \end{itemize}
    
    \note{About 45 seconds — reference slide
    ``Three questions, mirroring C2's structure.
    }
    \note[item]{Learning Efficiency}
    \note[item]{Robustness}
    \note[item]{Cost vs.\ Benefit}
\end{frame}


\begin{frame}[shrink=20]{Training Results}
    \vspace{1.1em}
    \begin{figure}
        \begin{center}
            \includegraphics[width=0.65\linewidth]{c3_training_curves.png}
            \caption{Evaluation performance across all conditions
                and sensor configurations.
                Implicit Indication (red) consistently
                outperforms PIC (green) and HAPPO (blue).}
            \label{fig:c3-eval-boxes}
        \end{center}
    \end{figure}

    \note{More of the same. For better or worse.}
\end{frame}


\begin{frame}{Evaluations}
    \begin{figure}
        \centering
        \includegraphics[width=\linewidth]{eval_heatmaps.png}
        \caption{Evaluation heatmaps: Implicit Indication (left),
                PIC (center), difference (right).
                Higher values (blue) in the difference map
                indicate Implicit Indication outperforms.}
        \label{fig:c3-heatmaps}
    \end{figure}
\end{frame}


\begin{frame}[shrink=7]{Interpreting the Gap}
    \vspace{1.5em}

    \textbf{PIC was designed for inter-agent structure:}
    \begin{itemize}
        \item The GNN critic models \emph{relationships between agents}:
            it excels when coordination patterns are rich
            and the graph topology is informative.
        \item In HyperGrid, heterogeneity lives at the
            \textbf{observation-element level} within each agent,
            not in the relational structure between agents.
    \end{itemize}

    \vspace{0.4em}
    \textbf{The signal problem:}
    \begin{itemize}
        \item We hoped PIC's graph structure over observation elements
            would help the critic learn \emph{which elements matter} —
            but the value signal it provides to the actors
            is too coarse to guide observation-conditioned behavior.
        \item Implicit Indication gives each actor \textbf{direct access}
            to observation availability via masking;
            no intermediary architecture required.
    \end{itemize}

    \vspace{0.3em}
    \begin{alertblock}{Takeaway}
        When heterogeneity is \emph{observational}, the actor needs
        direct information about what it can see; 
        not a critic-mediated summary of agent relationships.
    \end{alertblock}

    \note{About 1.5 minutes.
    Frame this charitably toward PIC:
    ``PIC is a good tool for the problem it was designed for.
    When agents interact richly — think formation control, communication
    networks — the graph structure carries real information, and a GNN
    critic can exploit that to provide strong value estimates.''
    Then the pivot:
    ``In HyperGrid, heterogeneity is not about how agents relate to
    each other. It is about which observation channels each agent has.
    We hoped the GNN, with observation elements as nodes within each agent,
    would learn to value those element-level distinctions.
    But the critic's signal passes through graph aggregation before
    reaching the actor — by the time it arrives, the fine-grained
    observation structure has been compressed away.''
    Then contrast:
    ``Implicit indication sidesteps this entirely. The actor sees the
    mask directly. Zero means `I cannot observe this.' The gradient
    for those input weights is exactly zero for that agent type.
    No intermediary, no compression, no information loss.''
    Deliver the alertblock to close.}
\end{frame}


% % -----------------------------------------------------------------------------
% % Slide 5.6 — C3 Summary
% % -----------------------------------------------------------------------------
% \begin{frame}{C3 Summary}

%     \vspace{-0.2em}
%     \begin{center}
%     \renewcommand{\arraystretch}{1.5}
%     \small
%     \begin{tabular}{@{} l p{8cm} @{}}
%         \toprule
%         \textbf{Dimension} & \textbf{Finding} \\
%         \midrule
%         Training
%             & Similar convergence dynamics; no clear training-time advantage
%               for either paradigm. \\
%         Evaluation
%             & Implicit Indication outperforms PIC by $2.2$--$3.7\times$
%               across all sensor configurations. \\
%         Robustness
%             & Implicit Indication dominates across all 8 perturbation conditions;
%               also outperforms the HAPPO baseline. \\
%         Cost
%             & Representational approach requires no specialized architectural
%               components — standard MLP/CNN suffices. \\
%         \bottomrule
%     \end{tabular}
%     \end{center}

%     \vspace{0.5em}
%     \begin{alertblock}{C3 Verdict}
%         In domains where heterogeneity stems from differing observation capabilities,
%         representational invariance via explicit masking outperforms
%         architectural invariance via graph neural networks —
%         with lower implementation complexity.
%     \end{alertblock}

%     \vspace{0.3em}
%     \centering
%     \small\emph{Three contributions, three results — all pointing the same direction.}

%     \note{About 1 minute.
%     Walk down the table briefly — one sentence per row.
%     ``Training: both learn at the same rate. Evaluation: implicit indication
%     wins by two to nearly four times. Robustness: implicit indication dominates
%     every perturbation. Cost: no GNN, no graph construction — just a
%     standard network with a well-designed input.''
%     Read the alertblock verdict.
%     Then deliver the transition line at the bottom:
%     ``Three contributions, three results — all pointing the same direction.
%     The next section pulls them together.''
%     This should feel like the climax resolving;
%     the synthesis section that follows is the denouement.}
% \end{frame}

